\chapter{Losses: L1, the spectral angle and the ramp}\label{ch:losses}
\begingroup\emergencystretch=3em % long unbreakable source tags; local to this chapter

\begin{onething}
Indradhanu is trained to minimise one number, $L = L_1 + \lambda(t)\,\SAM$: the masked mean
absolute error (how far each band is off, in reflectance) plus a weighted spectral angle (how
far the \emph{shape} of each spectrum is off, in radians), with the weight $\lambda(t)$ rising
from zero over the first epochs so that brightness is learned before shape.
\end{onething}

Every chapter so far has described what the network \emph{computes}. None has said how it comes
to compute anything useful. A freshly built Indradhanu has about five and a half million
parameters set to random values (\cref{ch:params}); its reconstructions are noise. Training is
the process that turns those random numbers into useful ones, and training needs a single
number that says how wrong the network currently is. That number is the \emph{loss}. This
chapter builds Indradhanu's loss piece by piece, and computes its derivatives by hand, because
the derivatives are what training actually uses (\cref{ch:trainloop}).

\section{What a loss is, and what it must do}\label{sec:losses:what}

\begin{definition}[Loss]\index{loss}
A \emph{loss} (or \emph{loss function}, or \emph{objective}) is a function that takes the
network's output and the truth it should have produced, and returns one non-negative real
number. Smaller means better. Training changes the network's parameters so as to make this
number smaller.
\end{definition}

Recall from \cref{ch:shapes} that a \emph{parameter}\index{parameter} is one adjustable number
inside the network (one entry of a weight matrix or a bias vector), and from \cref{ch:masking}
that the network is shown a patch with most of its tokens hidden and must reconstruct the
hidden part. Here the truth is the patch the satellite recorded, $\vect{x}$, and the output is
the reconstruction, $\xhat$. Both have shape \shape{(B, 165, 128, 128)}: a batch of $B$ patches,
165 bands, 128 $\times$ 128 pixels.

A loss must have two properties.

\begin{enumerate}
  \item \textbf{Smaller must mean better, always.} If two reconstructions differ and one is
  closer to the truth in the sense you care about, it must get the smaller loss. Otherwise
  training improves the wrong thing.
  \item \textbf{It must be differentiable} (almost everywhere). Training does not search
  blindly; it asks, for every parameter, ``if I nudge you up a little, does the loss go up or
  down, and how fast?'' That rate is a derivative. \Cref{sec:losses:chain} shows how it is
  computed through the whole network; this chapter computes it for the loss itself.
\end{enumerate}

The obvious first choice, ``count the pixels that are wrong'', fails property 2: a count does
not change when a reflectance moves from 0.2001 to 0.2002, so its derivative is zero almost
everywhere and gives no direction. The losses below all change smoothly (or at least
continuously) as $\xhat$ moves.

\begin{analogy}
The loss is the marking scheme of an exam the network sits millions of times. L1 marks every
band separately, like a spelling test: each wrong letter costs marks in proportion to how wrong
it is. The spectral angle marks the \emph{handwriting}: it ignores how large you wrote and asks
only whether the letters have the right shape. A good marking scheme for a hyperspectral
model needs both, and, like a teacher, it asks for legibility before style (the ramp of
\cref{sec:losses:ramp}).
\end{analogy}

\section{The masked mean absolute error}\label{sec:losses:l1}

\subsection{Which pixels are graded}

The loss is not computed over every pixel of the patch. A pixel is \emph{graded} only if all
three of the following hold.

\begin{enumerate}
  \item It was a \textbf{prediction target}: its token was hidden from the encoder
  (\cref{ch:masking}). Grading pixels the model could see would reward copying.
  \item It is \textbf{valid}: a real measurement, not nodata or cloud-filled.
  \item It is \textbf{interior}: not inside the erosion collar around an invalid region,
  where the overlapping patch embedding has mixed invalid neighbours into it (\cref{ch:masking}).
\end{enumerate}

Conditions 2 and 3 together are the \emph{eroded mask}; condition 1 is the pixel prediction
mask. Both are arrays of 0s and 1s of shape \shape{(B, 1, H, W)}, and the loss mask is their
product:
\begin{lstlisting}
pixel_pred_mask = self._pred_mask_to_pixel_mask(pred_mask, H, W)            # (B, 1, H, W)
eroded_mask     = TokenMasking.erode_mask(mask, kernel_size=cfg.erosion_kernel_size)
loss_mask       = pixel_pred_mask * eroded_mask                             # (B, 1, H, W)
\end{lstlisting}

\begin{definition}[Loss mask]\index{loss mask}
The \emph{loss mask} $M$ is the set of pixels $p$ whose mask value is 1, that is, pixels that
were hidden, valid and interior. Its size $\abs{M}$ is the number of graded pixels.
\end{definition}

Multiplying two 0/1 numbers is a logical AND: $1 \times 1 = 1$, and any product with a 0 in
it is 0. Four pixels show every case.

\begin{center}
\begin{tabular}{ccccl}
\toprule
pixel & hidden? & valid and interior? & product & graded? \\
\midrule
A & 1 & 1 & $1 \times 1 = 1$ & yes \\
B & 1 & 0 & $1 \times 0 = 0$ & no: nodata, or on a border \\
C & 0 & 1 & $0 \times 1 = 0$ & no: the model could see it \\
D & 0 & 0 & $0 \times 0 = 0$ & no \\
\bottomrule
\end{tabular}
\end{center}

\coursefig[0.95\linewidth]{iw-23-loss-mask}{Which pixels are graded. The pixel prediction mask
(hidden tokens, amber) is multiplied by the eroded validity mask (indigo; the cloud and its
collar hatched). Multiplying 0/1 masks is a logical AND, so only pixels that were hidden, valid
and away from invalid regions enter the loss. Right: the L1 term for one graded pixel.}{fig:losses:mask}

\subsection{From the token grid to the pixel grid}

The prediction mask is born on the token grid: one flag per stage-1 token, shape
\shape{(B, N)} with $N = 32 \times 32 = 1024$ for a 128-pixel patch (stage 1 has stride 4,
\cref{ch:patchembed}). The loss lives on the 128 $\times$ 128 pixel grid, so each flag must be
copied to the 4 $\times$ 4 = 16 pixels its token covers:
\begin{lstlisting}
token_grid = pred_mask.reshape(B, 1, H_tokens, W_tokens)                    # (B, 1, 32, 32)
return F.interpolate(token_grid, size=(H, W), mode="nearest")               # (B, 1, 128, 128)
\end{lstlisting}
\code{mode="nearest"} copies values; it never averages. Bilinear interpolation would produce
0.25, 0.5 and 0.75 at block edges, and the later test \code{loss_mask == 1} would silently
reject those pixels, leaving a lattice of ungraded gaps. Use nearest-neighbour for masks,
always.

\subsection{The definition, and the order of averaging}

Write $C$ for the number of bands (165), $e_{c,p} = \hat{x}_{c,p} - x_{c,p}$ for the
\emph{error} (or \emph{residual}) of band $c$ at pixel $p$, and $M$ for the graded pixels.

\begin{definition}[Masked mean absolute error]\index{L1 loss}\index{mean absolute error}
\begin{equation}\label{eq:losses:l1}
L_1 \;=\; \frac{1}{\abs{M}} \sum_{p \in M} \left( \frac{1}{C} \sum_{c=1}^{C} \abs{e_{c,p}} \right)
\;=\; \frac{1}{C\,\abs{M}} \sum_{p \in M} \sum_{c=1}^{C} \abs{\hat{x}_{c,p} - x_{c,p}} .
\end{equation}
\end{definition}

The name comes from the \emph{L1 norm}\index{L1 norm}, $\norm{\vect{v}}_1 = \sum_i \abs{v_i}$,
the sum of absolute values. Its sibling, the L2 norm $\norm{\vect{v}}_2 = \sqrt{\sum_i v_i^2}$,
gives the mean squared error when you square it and average.

The two forms in \cref{eq:losses:l1} are equal. To see it, move the constant $1/C$ out of the
inner bracket:
\begin{align*}
\frac{1}{\abs{M}} \sum_{p \in M} \frac{1}{C} \sum_{c} \abs{e_{c,p}}
&= \frac{1}{\abs{M}} \cdot \frac{1}{C} \sum_{p \in M} \sum_{c} \abs{e_{c,p}}
  && \text{($1/C$ is the same for every $p$)} \\
&= \frac{1}{C\,\abs{M}} \sum_{p \in M} \sum_{c} \abs{e_{c,p}} .
  && \text{(multiply the two fractions)}
\end{align*}
The step works only because every pixel has the same number of bands $C$. A ``mean of means''
equals the overall mean when the groups are the same size; it would not if some pixels had
fewer bands. (That is exactly the situation in the SpectralEarth trainer, where dead bands are
skipped per pixel; there each pixel divides by its own count of valid bands,
\cref{ch:pretrain}.)

The code follows the first form, band mean first:
\begin{lstlisting}
per_pixel_l1      = (x_hat - pixels).abs()                    # (B, 165, H, W)
per_pixel_l1_mean = per_pixel_l1.mean(dim=1, keepdim=True)    # (B, 1, H, W)  one value per pixel
valid_l1          = per_pixel_l1_mean[loss_mask == 1]         # (|M|,)  graded pixels only
l1_loss           = valid_l1.mean()                           # ()  a single number
\end{lstlisting}
Line 3 is \emph{boolean indexing}: it pulls the graded pixels out into a flat list, so the mean
on line 4 automatically divides by $\abs{M}$ and never counts a masked zero as if it were a
perfect prediction.

\begin{worked}[label={we:losses:l1}]{Masked L1 on three bands and two graded pixels}
Do this one on paper. Use three bands instead of 165, and three pixels, of which A and B are
graded and C was visible to the model (so its loss-mask value is 0).
\[
\begin{array}{lccc}
 & \text{band 1} & \text{band 2} & \text{band 3} \\
\vect{x}_A & 0.100 & 0.200 & 0.300 \\
\xhat_A & 0.120 & 0.180 & 0.330 \\[2pt]
\vect{x}_B & 0.050 & 0.100 & 0.080 \\
\xhat_B & 0.060 & 0.070 & 0.080 \\[2pt]
\vect{x}_C & 0.200 & 0.200 & 0.200 \\
\xhat_C & 0.500 & 0.500 & 0.500
\end{array}
\]
\textbf{Step 1, errors at A.}
\begin{align*}
e_{1,A} &= 0.120 - 0.100 = +0.020\\
e_{2,A} &= 0.180 - 0.200 = -0.020\\
e_{3,A} &= 0.330 - 0.300 = +0.030
\end{align*}
\textbf{Step 2, absolute values at A} (an error of $-0.02$ is as wrong as $+0.02$):
$0.020,\ 0.020,\ 0.030$.

\textbf{Step 3, band mean at A.}
\begin{align*}
0.020 + 0.020 + 0.030 &= 0.070\\
0.070 / 3 &= 0.023333
\end{align*}
\textbf{Step 4, the same at B.}
\begin{align*}
e_{1,B} &= 0.060 - 0.050 = +0.010\\
e_{2,B} &= 0.070 - 0.100 = -0.030\\
e_{3,B} &= 0.080 - 0.080 = \phantom{+}0.000\\
0.010 + 0.030 + 0.000 &= 0.040\\
0.040 / 3 &= 0.013333
\end{align*}
\textbf{Step 5, pixel C.} Its band mean is $0.300$, a terrible reconstruction, but its mask
value is 0, so it is not in $M$ and contributes nothing.

\textbf{Step 6, mean over the graded pixels}, $\abs{M} = 2$:
\begin{align*}
0.023333 + 0.013333 &= 0.036667\\
0.036667 / 2 &= 0.018333
\end{align*}
\textbf{Check with the second form} of \cref{eq:losses:l1}: the six absolute errors sum to
$0.070 + 0.040 = 0.110$, and $C\abs{M} = 3 \times 2 = 6$, so $0.110 / 6 = 0.018333$. The same.

The real code, run on these numbers, prints $L_1 = 0.018333$ \tagC{} (script
\file{research/textbook/tools/compute_23-losses.py}).
\end{worked}

\paragraph{Is 0.018 good?} Reflectance sits roughly between 0.02 and 0.30, so an average error
of 0.018 is large; you could see it in a plot. The training configuration's own baseline table
expects $L_1$ of about 0.08--0.10 at the first epoch and 0.003--0.008 after 500
epochs \tagM{configs/hyperspectral_segformer_exp_2.json}. So these made-up pixels are an early-training
reconstruction.

\subsection{The derivative of L1}\label{sec:losses:l1grad}

Training needs $\partial L_1 / \partial \hat{x}_{c,p}$ for every band and pixel: how fast the
loss changes when one predicted value moves.

First, the derivative of the absolute value. For $e > 0$, $\abs{e} = e$, whose derivative is
$+1$. For $e < 0$, $\abs{e} = -e$, whose derivative is $-1$. At $e = 0$ the graph has a corner
and no derivative exists; PyTorch uses $0$ there. All three cases together are the
\emph{sign function}\index{sign function}:
\[
\frac{d\abs{e}}{de} = \sign(e) = \begin{cases} +1 & e > 0\\ 0 & e = 0 \\ -1 & e < 0.\end{cases}
\]

Now differentiate \cref{eq:losses:l1} with respect to one predicted value $\hat{x}_{c,p}$ of a
graded pixel. Only one term of the double sum contains it, the term $\abs{\hat{x}_{c,p} -
x_{c,p}}$; every other term is a constant as far as $\hat{x}_{c,p}$ is concerned and has
derivative zero. The derivative of $e_{c,p} = \hat{x}_{c,p} - x_{c,p}$ with respect to
$\hat{x}_{c,p}$ is 1. So
\begin{equation}\label{eq:losses:l1grad}
\frac{\partial L_1}{\partial \hat{x}_{c,p}} =
\begin{cases}
\dfrac{\sign(e_{c,p})}{C\,\abs{M}} & p \in M,\\[8pt]
0 & p \notin M.
\end{cases}
\end{equation}
Write $n = C\abs{M}$ for the number of graded (band, pixel) entries. Every graded entry gets a
push of the same size, $1/n$; only its direction depends on the error.

\begin{worked}[label={we:losses:l1grad}]{The gradient of L1 for the two-pixel example}
In \cref{we:losses:l1}, $n = 3 \times 2 = 6$, so $1/n = 0.16667$. Read off the signs of the
errors:
\[
\begin{array}{lccc}
 & \text{band 1} & \text{band 2} & \text{band 3}\\
\text{signs at A} & + & - & + \\
\text{gradient at A} & +0.16667 & -0.16667 & +0.16667\\
\text{signs at B} & + & - & 0\\
\text{gradient at B} & +0.16667 & -0.16667 & 0\\
\text{gradient at C} & 0 & 0 & 0
\end{array}
\]
PyTorch's autograd on the same tensors returns exactly this table \tagC. Two things to notice.
Band 3 at B is already exact, so it gets no push. Pixel C, which has the largest errors of
all, gets no push either, because it was not graded.
\end{worked}

\section{Why L1 and not squared error}\label{sec:losses:l1l2}

The most common default loss is the \emph{mean squared error}\index{mean squared error} (MSE),
$\frac{1}{n}\sum e^2$, whose derivative is $\frac{2e}{n}$: proportional to the error. L1's
derivative is $\frac{\pm 1}{n}$: the same size however large the error. That one difference
decides what each loss learns.

\subsection{L1 learns a median, L2 learns a mean}

Strip the problem to its bones. You must predict one number $c$ for five observations $y_1,
\dots, y_5$, and you are graded by either
\[
S_2(c) = \sum_i (y_i - c)^2 \qquad\text{or}\qquad S_1(c) = \sum_i \abs{y_i - c}.
\]

\paragraph{The best $c$ under $S_2$ is the mean.} Differentiate $S_2$ with respect to $c$. By
the chain rule (stated properly in \cref{sec:losses:chain}; here you need only that the
derivative of $(y_i - c)^2$ is $2(y_i - c)\cdot(-1)$):
\begin{align*}
\frac{dS_2}{dc} &= \sum_i 2(y_i - c)(-1) = -2\sum_i (y_i - c).
\intertext{Set it to zero, the condition for the bottom of the bowl:}
0 &= \sum_i y_i - 5c \\
c &= \frac{1}{5}\sum_i y_i .
\end{align*}

\paragraph{The best $c$ under $S_1$ is the median.} Differentiate $S_1$:
\[
\frac{dS_1}{dc} = \sum_i \sign(c - y_i) = (\text{number of } y_i \text{ below } c) -
(\text{number of } y_i \text{ above } c).
\]
This is negative while more observations lie above $c$ than below (so moving $c$ up lowers
$S_1$), positive once more lie below, and it changes sign at the middle observation: the median.

\begin{worked}[label={we:losses:median}]{A glitch pulls the mean, not the median}
Take five reflectances of the same field, one of which is a sensor glitch:
$y = (0.10,\ 0.11,\ 0.12,\ 0.13,\ 0.90)$.

\textbf{Mean.}
\begin{align*}
0.10 + 0.11 + 0.12 + 0.13 + 0.90 &= 1.36\\
1.36 / 5 &= 0.272
\end{align*}
\textbf{Median.} Sorted, the middle value is $0.12$.

The mean, $0.272$, is larger than four of the five observations; it describes none of them. The
median, $0.12$, describes the field.

\textbf{The pull of each observation at $c = 0.12$.} Under $S_2$ each observation pulls with
$2(c - y_i)$:
\begin{align*}
2(0.12 - 0.10) &= +0.04\\
2(0.12 - 0.11) &= +0.02\\
2(0.12 - 0.12) &= \phantom{+}0.00\\
2(0.12 - 0.13) &= -0.02\\
2(0.12 - 0.90) &= -1.56
\end{align*}
The glitch pulls 39 times harder than the next largest term ($1.56 / 0.04 = 39$). Under $S_1$
each pull is $\sign(c - y_i)$: $+1, +1, 0, -1, -1$. The glitch is one vote among five.
\end{worked}

\subsection{Why that matters for this training set}

The training corpus is real satellite imagery, not curated. It contains genuine anomalies (the
very things the model should find surprising), dead detectors and stripes, and cloud edges
the masks missed. Under MSE, those few pixels with huge errors would dominate every gradient,
and the model would spend its capacity learning to reconstruct exactly what it should fail on.
Under L1, each pixel's pull has the same size; outliers vote, but they do not take over. The
model learns the bulk of the data, which is normality, and that is what an anomaly detector
needs (\cref{ch:recon}).

\begin{center}
\begin{tabular}{lll}
\toprule
 & L1 & MSE \\
\midrule
penalty for error $e$ & $\abs{e}$ & $e^2$\\
a 2 $\times$ larger error costs & 2 $\times$ & 4 $\times$\\
a 10 $\times$ larger error costs & 10 $\times$ & 100 $\times$\\
pull (derivative) of one entry & $\pm 1/n$ & $2e/n$\\
best constant prediction & median & mean\\
dominated by outliers? & no & yes\\
\bottomrule
\end{tabular}
\end{center}

The term for a loss whose pull is bounded is \emph{robust}\index{robust loss}; the classical
reference for robust estimation of a location is \citet{huber1964robust}.

\begin{sourcenote}[: Antardhana and Tirohita are not identical apart from the loss]
The course (IW23) says the repo holds two models that differ \emph{only} in loss, Antardhana
(MSE) and Tirohita (L1). The model register \file{research/model_break_down/README.md} (rows 3
and 7) lists them as the same family, the spatial masked autoencoder, but with different
parameter counts, 330,314 for Tirohita and 265,480 for
Antardhana \tagM{research/model_break_down/README.md}, and different versions. So they are
siblings with the loss as the headline difference, not a controlled one-variable comparison.
The robustness argument above stands on its own; treat the pair as a hint, not as proof.
\end{sourcenote}

\section{The 40\% rule and sample counting}\label{sec:losses:forty}

Before any loss is computed, whole patches are thrown away:
\begin{lstlisting}
MIN_VALID_PIXEL_FRACTION = 0.4
valid_fractions = mask.flatten(1).float().mean(dim=1)     # (B,)  mean of 0/1 = fraction valid
keep = valid_fractions >= MIN_VALID_PIXEL_FRACTION        # (B,)  bool
num_kept = keep.sum().item()                              # int
\end{lstlisting}
(\file{app/foundation_models/trainers/hyperspectral_segformer_mae_trainer.py}, lines 34 and
162--170.) A patch less than 40\% valid is mostly nodata: training on it teaches the model to
reconstruct holes, and its few real pixels all lie near a border. The mean of a 0/1 array is
the fraction of 1s, the same trick as in \cref{ch:masking}.

\code{compute_loss} returns the loss together with \code{num_kept}, and the training loop
weights by it (\cref{ch:trainloop}). If no patch survives, the loss is a zero tensor and the
loop skips the batch. The function also returns zero when the loss mask is empty, so a
batch with nothing graded cannot divide by zero.

\begin{sourcenote}[: training and inference use different thresholds]
Training drops patches below 40\% valid. The scene-scale inference loop in
\file{app/foundation_models/inferencers/segformer_mae_inferencer.py} (lines 281--284), which the
hyperspectral inferencer inherits, uses \code{MIN_VALID_FRACTION = 0.1} under a comment that
says ``Matches training: patches with $<$ 40\% valid pixels were discarded.'' The comment is
wrong; the constant is 0.1. The looser inference threshold is defensible (you would rather
have an imperfect score near a scene edge than none), but if you debug edge artefacts, know
that these two numbers differ. \Cref{ch:inference} returns to this.
\end{sourcenote}

\section{The spectral angle}\label{sec:losses:sam}

\subsection{Why L1 is not enough}

L1 measures error in reflectance units, and so it is dominated by bright pixels. A bright pixel
at reflectance 0.30 that is 10\% off costs $0.03$ per band; a dark pixel at 0.03 that is 10\%
off costs $0.003$. Same proportional error, ten times the penalty, so gradient descent spends
its effort on bright surfaces and neglects shadow, dark water and dark soil.

Yet the diagnostic content of a spectrum is its \emph{shape}, not its level
(\cref{ch:light}). Sunlit grass and shadowed grass are the same material with the same shape at
different brightness; a gas plume changes the shape at a few wavelengths without changing the
overall level much. You need a second measurement that sees shape and ignores brightness.

\subsection{The definition}

Treat each pixel's $C$ values as a vector in $\R^C$ (\cref{ch:linalg}). The angle between two
vectors measures how differently they point, regardless of length.

\begin{definition}[Spectral angle]\index{spectral angle}\index{SAM|see{spectral angle}}
For nonzero spectra $\vect{x}, \xhat \in \R^C$,
\begin{equation}\label{eq:losses:sam}
u = \cos\theta = \frac{\ip{\vect{x}}{\xhat}}{\norm{\vect{x}}\,\norm{\xhat}},
\qquad
\SAM(\vect{x}, \xhat) = \theta = \arccos u \in [0, \pi],
\end{equation}
where $\ip{\vect{x}}{\xhat} = \sum_c x_c \hat{x}_c$ is the dot product and $\norm{\vect{x}} =
\sqrt{\ip{\vect{x}}{\vect{x}}}$ the length. The angle is in radians.
\end{definition}

The name is \emph{Spectral Angle Mapper}\index{Spectral Angle Mapper}, from the remote-sensing
classifier of \citet{kruse1993spectral}, which matches an unknown spectrum against a library by
this angle. (The companion volume on spectral matching uses it that way; \cref{ch:material}
uses it at the end of the chain.) Here it is repurposed as a training loss between the
observed and the reconstructed spectrum.

That $u$ really is the cosine of an angle, and lies in $[-1, 1]$, is the Cauchy--Schwarz
inequality $\abs{\ip{\vect{x}}{\xhat}} \le \norm{\vect{x}}\norm{\xhat}$ (\cref{ch:linalg}).

\begin{center}
\begin{tabular}{ll}
\toprule
angle & meaning \\
\midrule
0 & identical shape, whatever the brightness\\
0.01 to 0.10 rad (0.6\textdegree{} to 5.7\textdegree) & a good reconstruction here\\
$\pi/2 \approx 1.571$ & at right angles: unrelated\\
$\pi \approx 3.142$ & exactly opposite\\
\bottomrule
\end{tabular}
\end{center}

\begin{worked}[label={we:losses:sam}]{The spectral angle of pixel A, by hand}
Do this on paper. Pixel A of \cref{we:losses:l1}: $\vect{x} = (0.100, 0.200, 0.300)$, $\xhat =
(0.120, 0.180, 0.330)$.

\textbf{Step 1, the dot product.}
\begin{align*}
0.100 \times 0.120 &= 0.0120\\
0.200 \times 0.180 &= 0.0360\\
0.300 \times 0.330 &= 0.0990\\
0.0120 + 0.0360 + 0.0990 &= 0.1470
\end{align*}
\textbf{Step 2, the two lengths.}
\begin{align*}
\norm{\vect{x}}^2 &= 0.100^2 + 0.200^2 + 0.300^2 = 0.0100 + 0.0400 + 0.0900 = 0.1400\\
\norm{\vect{x}} &= \sqrt{0.1400} = 0.374166\\
\norm{\xhat}^2 &= 0.120^2 + 0.180^2 + 0.330^2 = 0.0144 + 0.0324 + 0.1089 = 0.1557\\
\norm{\xhat} &= \sqrt{0.1557} = 0.394588
\end{align*}
\textbf{Step 3, the cosine.}
\begin{align*}
\norm{\vect{x}}\,\norm{\xhat} &= 0.374166 \times 0.394588 = 0.147641\\
u &= 0.1470 / 0.147641 = 0.995655
\end{align*}
\textbf{Step 4, the angle.}
\begin{align*}
\theta &= \arccos(0.995655) = 0.09325 \text{ rad}\\
\theta &= 0.09325 \times 180 / \pi = 5.34\text{\textdegree}
\end{align*}
So pixel A has $L_1$ contribution $0.0233$ (its brightness is off by 0.023 reflectance on
average) and spectral angle $0.0933$ rad (its shape is off by about 5.3\textdegree). They
measure different things; neither implies the other.
\end{worked}

\begin{sourcenote}[: the course's angle for this pixel]
The course (IW24) gets $\norm{\vect{x}}\norm{\xhat} = 0.147644$, $u = 0.995637$ and $\theta =
0.0935$ rad $= 5.35$\textdegree. The product is $0.147641$ (more digits: $0.1476415$), which
gives $u = 0.995655$ and $\theta = 0.09325$ rad $= 5.34$\textdegree. A float64 evaluation and
the real \code{SAMLoss} (in float32, eps included) both give $0.09325$ \tagC{}
(\file{research/textbook/tools/compute_23-losses_output.txt}). The difference is a rounding slip
in the course's third step; it changes no conclusion. The same slip propagates into the
course's atan2 check, which gives $\sin$-term $0.013784$ instead of $0.013748$.
\end{sourcenote}

\subsection{Scale invariance}\label{sec:losses:scaleinv}

\begin{worked}[label={we:losses:double}]{A reconstruction twice as bright}
Suppose the model predicted $\xhat = (0.200, 0.400, 0.600)$, exactly twice the truth in
every band.

\textbf{L1.}
\begin{align*}
\abs{0.200 - 0.100} &= 0.100\\
\abs{0.400 - 0.200} &= 0.200\\
\abs{0.600 - 0.300} &= 0.300\\
(0.100 + 0.200 + 0.300)/3 &= 0.600/3 = 0.200
\end{align*}
Enormous: almost nine times pixel A's 0.0233.

\textbf{Spectral angle.}
\begin{align*}
\ip{\vect{x}}{\xhat} &= 0.100 \times 0.200 + 0.200 \times 0.400 + 0.300 \times 0.600 = 0.02 + 0.08 + 0.18 = 0.28\\
\norm{\xhat} &= \sqrt{0.04 + 0.16 + 0.36} = \sqrt{0.56} = 0.748331\\
\norm{\vect{x}}\,\norm{\xhat} &= 0.374166 \times 0.748331 = 0.28\\
u &= 0.28 / 0.28 = 1, \qquad \theta = \arccos 1 = 0.
\end{align*}
The angle reports a perfect match.
\end{worked}

This is not a coincidence of the numbers. Here is the general statement and its proof.

\paragraph{Claim.} For any $\alpha > 0$, $\SAM(\vect{x}, \alpha\xhat) = \SAM(\vect{x}, \xhat)$.

\paragraph{Proof.} Compute the cosine for $\alpha\xhat$ and simplify.
\begin{align*}
\frac{\ip{\vect{x}}{\alpha\xhat}}{\norm{\vect{x}}\,\norm{\alpha\xhat}}
&= \frac{\alpha\ip{\vect{x}}{\xhat}}{\norm{\vect{x}}\,\norm{\alpha\xhat}}
  && \text{(a constant factor comes out of a dot product)}\\
&= \frac{\alpha\ip{\vect{x}}{\xhat}}{\norm{\vect{x}}\,\sqrt{\textstyle\sum_c \alpha^2\hat{x}_c^2}}
  && \text{(definition of the length)}\\
&= \frac{\alpha\ip{\vect{x}}{\xhat}}{\norm{\vect{x}}\,\abs{\alpha}\,\norm{\xhat}}
  && \text{($\sqrt{\alpha^2} = \abs{\alpha}$)}\\
&= \frac{\ip{\vect{x}}{\xhat}}{\norm{\vect{x}}\,\norm{\xhat}}
  && \text{($\alpha > 0$, so $\abs{\alpha} = \alpha$ cancels).}
\end{align*}
The cosine is unchanged, so the angle is unchanged. $\square$

If $\alpha < 0$, $\abs{\alpha} = -\alpha$ and the cosine changes sign: the angle becomes
$\pi - \theta$. Reflectances are non-negative, so this case does not arise in practice.

\coursefig[0.95\linewidth]{iw-24-l1-sam}{Two complementary errors. (a) Treating a spectrum as a
vector, the spectral angle is the angle between reconstruction and truth, and L1 the average
per-band distance. A reconstruction that is the truth scaled by 2 (amber) has an angle of
exactly 0 but a large L1; a reconstruction close in level but slightly rotated (indigo) has a
small L1 and a small angle. (b) The three-band numbers of \cref{we:losses:sam}.}{fig:losses:l1sam}

So the model needs both terms:
\begin{center}
\begin{tabular}{lcc}
\toprule
kind of error & caught by L1 & caught by the angle\\
\midrule
brightness only (right shape, wrong scale) & yes & no\\
shape only (right level, wrong colour) & weakly & yes\\
both & yes & yes\\
\bottomrule
\end{tabular}
\end{center}

\section{The chain rule}\label{sec:losses:chain}

To train with the spectral angle you need its derivative with respect to every predicted
value, and the angle is a function of a function of a function of $\xhat$. The tool for that
is the chain rule. It is used everywhere in the rest of this book, so it is stated carefully
here and shown on a tiny example first.

\begin{definition}[Chain rule, one variable]\index{chain rule}
If $y = f(g(w))$, write $z = g(w)$. Then
\[
\frac{dy}{dw} = \frac{dy}{dz}\cdot\frac{dz}{dw} = f'(g(w))\, g'(w).
\]
In words: the rate at which the outer function changes with the input is the rate at which it
changes with its own input, times the rate at which that input changes with $w$.
\end{definition}

\begin{worked}[label={we:losses:chain}]{The chain rule on a tiny function}
Let $y = (3w + 1)^2$ and find $dy/dw$ at $w = 1$.

\textbf{Name the inner function.} $z = g(w) = 3w + 1$, so $y = f(z) = z^2$.

\textbf{Differentiate each piece.} $dz/dw = 3$, and $dy/dz = 2z$.

\textbf{Evaluate at $w = 1$.} $z = 3 \times 1 + 1 = 4$, so $dy/dz = 2 \times 4 = 8$.

\textbf{Multiply.} $dy/dw = 8 \times 3 = 24$.

\textbf{Check numerically.} Move $w$ a little either side and divide the change in $y$ by the
change in $w$:
\begin{align*}
w = 1.001: \quad y &= (3.003 + 1)^2 = 4.003^2 = 16.024009\\
w = 0.999: \quad y &= (2.997 + 1)^2 = 3.997^2 = 15.976009\\
\frac{16.024009 - 15.976009}{1.001 - 0.999} &= \frac{0.048}{0.002} = 24.
\end{align*}
The chain rule agrees with the measured slope.
\end{worked}

With several variables the rule is the same, summed over every path. If $y$ depends on $w$
through intermediate quantities $z_1, \dots, z_k$, then
\[
\frac{\partial y}{\partial w} = \sum_{j=1}^{k} \frac{\partial y}{\partial z_j}\,\frac{\partial z_j}{\partial w}.
\]
Training a network is this rule applied from the loss backwards through every layer, which is
why the procedure is called \emph{backpropagation}\index{backpropagation}
(\cref{ch:trainloop}).

Two more rules are used below.

\paragraph{The quotient rule.}\index{quotient rule} For $q = a/b$ where $a$ and $b$ both depend
on $w$: $\dfrac{dq}{dw} = \dfrac{a'b - ab'}{b^2}$. Tiny check: $q = w/(w+1)$ at $w = 1$ has $a'
= 1$, $b' = 1$, so $q' = (1 \cdot 2 - 1 \cdot 1)/2^2 = 1/4$. Numerically, $q(1.001) =
1.001/2.001 = 0.500250$ and $q(0.999) = 0.999/1.999 = 0.499750$, so the slope is $0.000500/0.002
= 0.25$.

\paragraph{The derivative of a length.} For $b = \norm{\xhat} = \big(\sum_c
\hat{x}_c^2\big)^{1/2}$, the chain rule with outer function $(\cdot)^{1/2}$ gives
\[
\frac{\partial b}{\partial \hat{x}_c} = \frac{1}{2}\Big(\sum_k \hat{x}_k^2\Big)^{-1/2}\cdot 2\hat{x}_c = \frac{\hat{x}_c}{\norm{\xhat}}.
\]
The gradient of a length is the unit vector pointing along it.

\section{The derivative of arccos, and where it blows up}\label{sec:losses:arccos}

\paragraph{Derivation.} Let $y = \arccos u$, so $\cos y = u$, with $y \in [0, \pi]$.
\begin{align*}
\cos y &= u
  && \text{(definition of arccos)}\\
-\sin y \,\frac{dy}{du} &= 1
  && \text{(differentiate both sides with respect to $u$; chain rule on the left)}\\
\frac{dy}{du} &= -\frac{1}{\sin y}
  && \text{(divide by $-\sin y$)}\\
&= -\frac{1}{\sqrt{1 - \cos^2 y}}
  && \text{($\sin^2 + \cos^2 = 1$, and $\sin y \ge 0$ on $[0, \pi]$)}\\
&= -\frac{1}{\sqrt{1 - u^2}} .
  && \text{(substitute $\cos y = u$)}
\end{align*}

\paragraph{Numbers.} Here is the size of that derivative as $u$ approaches 1, with the angle
that each $u$ corresponds to:
\begin{center}
\begin{tabular}{rrrrr}
\toprule
$u$ & $1 - u^2$ & $\sqrt{1 - u^2}$ & $d\arccos/du$ & angle \\
\midrule
0.9 & 0.19 & 0.4359 & $-2.29$ & 0.451 rad (25.8\textdegree)\\
0.99 & 0.0199 & 0.1411 & $-7.09$ & 0.142 rad (8.1\textdegree)\\
0.999 & 0.001999 & 0.04471 & $-22.4$ & 0.0447 rad (2.6\textdegree)\\
0.9999 & 0.00019999 & 0.01414 & $-70.7$ & 0.0141 rad (0.8\textdegree)\\
1 & 0 & 0 & $-\infty$ & 0\\
\bottomrule
\end{tabular}
\end{center}
Work the third row on paper: $0.999^2 = 0.998001$; $1 - 0.998001 = 0.001999$; $\sqrt{0.001999}
= 0.04471$; $1 / 0.04471 = 22.37$.

Now set beside this table where a \emph{good} model lives. The configuration's baseline
expects an angle of 0.02--0.05 rad after long training \tagM{configs/hyperspectral_segformer_exp_2.json},
that is $u$ between $\cos 0.05 = 0.99875$ and $\cos 0.02 = 0.99980$. A well-trained model sits
at $u = 0.999$ and above, in the bottom rows, where the factor $d\arccos/du$ is in the tens and
heading for infinity. The better the model gets, the closer it sits to the singularity.

\begin{figure}[htbp]
\centering
\begin{tikzpicture}
\begin{groupplot}[group style={group size=2 by 1, horizontal sep=1.7cm},
  width=0.47\linewidth, height=5.6cm, ymode=log, ymin=0.8, ymax=200,
  ylabel={$\abs{d\arccos/du} = 1/\sqrt{1-u^2}$}, grid=major, grid style={gray!20}]
\nextgroupplot[xlabel={$u = \cos\theta$}, xmin=0, xmax=1, title={(a) linear $u$}]
  % parametric in the angle (degrees): u = cos(a), y = 1/sin(a)
  \addplot[loss, thick, domain=0.4:90, samples=400, variable=\a] ({cos(\a)}, {1/sin(\a)});
  \addplot[only marks, mark=*, mark size=1.6pt, spat] coordinates
    {(0.9,2.294) (0.99,7.089) (0.999,22.37) (0.9999,70.71)};
  \draw[dashed, maskgrey] (axis cs:0.999,0.8) -- (axis cs:0.999,200);
  \node[font=\scriptsize, anchor=east, align=right] at (axis cs:0.985,60) {well-trained\\ $u \ge 0.999$};
\nextgroupplot[xlabel={$1 - u$ (log scale)}, xmode=log, xmin=1e-5, xmax=1, x dir=reverse,
  title={(b) zoom on $u \to 1$}]
  \addplot[loss, thick, domain=-5:0, samples=200, variable=\k] ({10^\k}, {1/sqrt((10^\k)*(2-10^\k))});
  \addplot[only marks, mark=*, mark size=1.6pt, spat] coordinates
    {(0.1,2.294) (0.01,7.089) (0.001,22.37) (0.0001,70.71)};
  \fill[spec, opacity=0.12] (axis cs:1e-5,0.8) rectangle (axis cs:1e-3,200);
  \node[font=\scriptsize, align=center] at (axis cs:1e-4,1.6) {good model\\ lives here};
\end{groupplot}
\end{tikzpicture}
\caption{The size of the derivative of arccos, drawn to scale from $1/\sqrt{1-u^2}$, on a
logarithmic vertical axis. Dots are the four table rows. (a) Against $u$ itself: flat until
about $u = 0.9$, then a wall. (b) Against $1 - u$ on a log axis, reading right to left as $u
\to 1$: a straight line of slope $-\frac12$, because $1/\sqrt{1-u^2} \approx 1/\sqrt{2(1-u)}$
there. The shaded band, $u \ge 0.999$, is where a well-trained model's pixels sit.}
\label{fig:losses:arccos}
\end{figure}

\section{The gradient of the spectral angle}\label{sec:losses:samgrad}

Now combine the pieces. Write $a = \norm{\vect{x}}$ (a constant, since the truth does not
depend on the parameters), $b = \norm{\xhat}$, $s = \ip{\vect{x}}{\xhat}$, so $u = s/(ab)$ and
$\theta = \arccos u$.

\paragraph{Step 1: the derivative of the cosine.} Use the quotient rule with numerator $s$ and
denominator $ab$. The pieces: $\partial s/\partial\hat{x}_c = x_c$ (only one term of the dot
product contains $\hat{x}_c$), and $\partial(ab)/\partial \hat{x}_c = a\,\hat{x}_c/b$ (the
derivative of a length, times the constant $a$).
\begin{align}
\frac{\partial u}{\partial \hat{x}_c}
&= \frac{x_c \cdot ab - s \cdot a\hat{x}_c/b}{(ab)^2}
  && \text{(quotient rule)} \notag\\
&= \frac{x_c}{ab} - \frac{s\,\hat{x}_c}{a b^3}
  && \text{(split the fraction; cancel $ab$ in the first, $a$ in the second)} \notag\\
&= \frac{1}{b}\left(\frac{x_c}{a} - u\,\frac{\hat{x}_c}{b}\right).
  && \text{(factor out $1/b$; $s/(ab) = u$)}\label{eq:losses:du}
\end{align}

\paragraph{Step 2: the chain rule through arccos.}
\begin{equation}\label{eq:losses:dtheta}
\frac{\partial\theta}{\partial \hat{x}_c} = \frac{d\theta}{du}\cdot\frac{\partial u}{\partial \hat{x}_c}
= -\frac{1}{\sqrt{1 - u^2}}\cdot\frac{1}{b}\left(\frac{x_c}{a} - u\,\frac{\hat{x}_c}{b}\right).
\end{equation}

\paragraph{Reading the formula.} The bracket is the unit truth $\vect{x}/a$ minus $u$ times the
unit reconstruction $\xhat/b$: the part of the truth's direction that is perpendicular to the
reconstruction. Two facts follow.

\begin{enumerate}
\item \textbf{The gradient is perpendicular to $\xhat$.} Take the dot product of the bracket
with $\xhat$: $\ip{\vect{x}}{\xhat}/a - u\,\ip{\xhat}{\xhat}/b = s/a - u b = ub - ub = 0$. So
the angle's gradient never asks the reconstruction to get brighter or darker, only to rotate.
That is scale invariance seen through the derivative.
\item \textbf{The product is finite, even though one factor is not.} The squared length of the
bracket is
\[
\norm{\tfrac{\vect{x}}{a} - u\tfrac{\xhat}{b}}^2 = 1 - 2u\cdot u + u^2 = 1 - u^2,
\]
(expand the square; the two unit vectors have length 1 and dot product $u$). So
$\norm{\partial u/\partial \xhat} = \sqrt{1 - u^2}/b$, and
\[
\norm{\frac{\partial\theta}{\partial\xhat}} = \frac{1}{\sqrt{1 - u^2}}\cdot\frac{\sqrt{1-u^2}}{b} = \frac{1}{\norm{\xhat}} .
\]
The $\sqrt{1 - u^2}$ blowing up in the arccos factor is cancelled exactly by the same quantity
shrinking in the cosine's gradient. In exact arithmetic the spectral angle's gradient has the
fixed length $1/\norm{\xhat}$ at every angle $\theta > 0$. Only at $\theta = 0$ is its direction
undefined: the angle has a corner there, as $\abs{e}$ has at $e = 0$.
\end{enumerate}

\begin{worked}[label={we:losses:samgrad}]{The gradient of the angle for pixel A}
Continue \cref{we:losses:sam}: $a = 0.374166$, $b = 0.394588$, $s = 0.147$, $u = 0.995655$.

\textbf{The first term of \cref{eq:losses:du}}, $x_c/(ab)$ with $ab = 0.147641$:
\begin{align*}
0.100/0.147641 &= 0.677317\\
0.200/0.147641 &= 1.354633\\
0.300/0.147641 &= 2.031950
\end{align*}
\textbf{The second term}, $s\,\hat{x}_c/(ab^3)$. First $b^3 = 0.394588^3 = 0.061437$, then
$ab^3 = 0.374166 \times 0.061437 = 0.022988$, and $s/(ab^3) = 0.147/0.022988 = 6.39470$:
\begin{align*}
6.39470 \times 0.120 &= 0.767364\\
6.39470 \times 0.180 &= 1.151047\\
6.39470 \times 0.330 &= 2.110252
\end{align*}
\textbf{Subtract} to get $\partial u/\partial\xhat$:
\begin{align*}
0.677317 - 0.767364 &= -0.090048\\
1.354633 - 1.151047 &= +0.203586\\
2.031950 - 2.110252 &= -0.078302
\end{align*}
\textbf{The arccos factor.} $1 - u^2 = 1 - 0.991329 = 0.008671$; $\sqrt{0.008671} = 0.093116$;
$-1/0.093116 = -10.7393$.

\textbf{Multiply} to get $\partial\theta/\partial\xhat$:
\begin{align*}
-10.7393 \times (-0.090048) &= +0.96705\\
-10.7393 \times (+0.203586) &= -2.18638\\
-10.7393 \times (-0.078302) &= +0.84092
\end{align*}
\textbf{Interpret.} Gradient descent moves $\xhat$ \emph{against} the gradient: down in
band 1, up in band 2, down in band 3. Compare the unit vectors: $\vect{x}/a = (0.267, 0.535,
0.802)$ and $\xhat/b = (0.304, 0.456, 0.836)$. The reconstruction leans too far towards bands
1 and 3 and not enough towards band 2; the gradient corrects exactly that.

\textbf{Check the two facts.} The dot product with $\xhat$:
$0.120(-0.090048) + 0.180(0.203586) + 0.330(-0.078302) = -0.010806 + 0.036645 - 0.025840 =
-0.000001 \approx 0$. The length: $\sqrt{0.96705^2 + 2.18638^2 + 0.84092^2} = \sqrt{0.93519 +
4.78026 + 0.70715} = \sqrt{6.42260} = 2.5343$, and $1/b = 1/0.394588 = 2.5343$.

PyTorch's autograd through \code{torch.acos} in float64 returns $(0.96705, -2.18638,
0.84092)$ \tagC; the real \code{SAMLoss} in float32 returns $(0.96701, -2.18630,
0.84087)$ \tagC{} (\file{research/textbook/tools/compute_23-losses_output.txt}).
\end{worked}

\section{Why the code does not call arccos}\label{sec:losses:atan2}

If the gradient's length is a tame $1/\norm{\xhat}$, what is the danger? The danger is in how a
computer evaluates \cref{eq:losses:dtheta}. Automatic differentiation computes the two factors
separately and then multiplies them. The arccos factor is evaluated from $u$ as computed in
floating point, and close to $u = 1$ floating point cannot tell nearby angles apart.

\paragraph{How coarse is float32 near 1?} A float32 number has 24 significant bits, so the gap
between 1 and the next smaller representable number is $2^{-24} = 5.96 \times 10^{-8}$. For a
small angle, $\cos\theta \approx 1 - \theta^2/2$. When $\theta^2/2$ is less than half the gap,
$2.98 \times 10^{-8}$, the cosine rounds to exactly 1:
\[
\frac{\theta^2}{2} < 2.98\times 10^{-8} \iff \theta < \sqrt{5.96\times 10^{-8}} = 2.44 \times 10^{-4}\text{ rad}.
\]
Indeed, in float32 $\cos(3\times 10^{-4}) = 0.99999994$ but $\cos(2 \times 10^{-4}) = 1.0$
exactly \tagC. Below about $2.4 \times 10^{-4}$ rad (0.014\textdegree) the cosine carries no
information, and the computed $u$ is rounding noise that can even exceed 1.

\begin{worked}[label={we:losses:nan}]{The naive formula producing NaN}
Feed a perfect reconstruction, $\xhat = \vect{x}$, through
\code{torch.acos(dot / (norm_x * norm_hat))} in float32 and ask for the gradient
(\file{research/textbook/tools/compute_23-losses.py}).
\begin{itemize}
\item $\vect{x} = (0.2, 0.2, 0.2)$: the computed $u$ is $1.0000001$, above 1. arccos of it is
NaN (``not a number''), and so is the gradient. With the usual guard \code{u.clamp(-1, 1)} the
angle becomes 0 and the gradient becomes exactly 0 (a clamp passes no gradient once it is
active), so training silently stops pushing.
\item $\vect{x} = (0.25, 0.5, 0.125)$: $u$ comes out exactly $1.0$. arccos gives 0, but the
factor $-1/\sqrt{1 - u^2}$ is $-\infty$ and $\partial u/\partial\xhat$ is 0, and $\infty \times 0$
is NaN. The clamp does not help here.
\end{itemize}
One NaN in one pixel's gradient spreads through the optimiser into every parameter at the
next step. And it appears precisely when the model has become good: late in a long, expensive
run.
\end{worked}

\subsection{The atan2 formulation}

\code{SAMLoss} (\file{app/foundation_models/components/sam_loss.py}) rewrites the angle with a
sine as well as a cosine. With $N = \norm{\vect{x}}\norm{\xhat}$:
\begin{align*}
N\cos\theta &= \ip{\vect{x}}{\xhat} && \text{(the dot product)}\\
N\sin\theta &= \sqrt{N^2 - \ip{\vect{x}}{\xhat}^2} && \text{(from $\sin^2\theta = 1 - \cos^2\theta$, times $N^2$)}\\
\theta &= \operatorname{atan2}(N\sin\theta,\; N\cos\theta),
\end{align*}
where $\operatorname{atan2}(y, x)$ is the angle of the point $(x, y)$ measured from the positive
$x$-axis, in $[0, \pi]$ when $y \ge 0$. Its derivatives, $\partial/\partial y = x/(x^2 + y^2)$
and $\partial/\partial x = -y/(x^2 + y^2)$, are finite everywhere except at the origin.

The code, with its guards:
\begin{lstlisting}
x_m  = x * mask                                                        # (B, 165, H, W)
xh_m = x_hat * mask
dot      = (xh_m * x_m).sum(dim=1, keepdim=True)                       # (B, 1, H, W)
norm_hat = (xh_m * xh_m).sum(dim=1, keepdim=True).add(self.eps).sqrt() # guard 0
norm_x   = (x_m * x_m).sum(dim=1, keepdim=True).add(self.eps).sqrt()
cross_norm  = (norm_hat * norm_x).clamp(min=self.eps)                  # guard 1
cos_term    = dot.clamp(-cross_norm, cross_norm)                       # guard 2
sin_term_sq = (cross_norm ** 2 - cos_term ** 2).clamp(min=0)
sin_term    = (sin_term_sq + self.eps).sqrt()                          # guard 3
angles      = torch.atan2(sin_term, cos_term)                          # (B, 1, H, W), in [0, pi]
\end{lstlisting}
with $\varepsilon = 10^{-8}$ (\code{eps}, line 30).

\begin{description}[style=nextline,leftmargin=1.2em]
\item[Guard 0, $+\varepsilon$ inside each length.] A masked pixel is a zero vector; its length
would be $\sqrt{0}$, whose derivative $1/(2\sqrt{0})$ is infinite.
\item[Guard 1, \texttt{clamp(min=eps)} on the product of lengths.] Prevents a zero (or
vanishing) $N$.
\item[{Guard 2, clamp the dot product to $[-N, N]$.}] Mathematically $\abs{\ip{\vect{x}}{\xhat}}
\le N$ (Cauchy--Schwarz); in floating point it can exceed it by a hair, as in
\cref{we:losses:nan}, and then $N^2 - \text{dot}^2$ would be negative.
\item[Guard 3, $+\varepsilon$ inside the sine's square root.] For a perfect reconstruction
$N^2 - \text{dot}^2 = 0$, and $d\sqrt{v}/dv = 1/(2\sqrt{v})$ is infinite at $v = 0$: the same
disease as arccos, one level down. Adding $\varepsilon$ keeps the sine term at least
$\sqrt{10^{-8}} = 10^{-4}$, so the derivative of the square root is at most $1/(2 \times
10^{-4}) = 5000$.
\end{description}

\begin{sourcenote}[: the module docstring, and the guard the course leaves out]
The module docstring of \file{app/foundation_models/components/sam_loss.py} (line 9) still
writes the loss as $\arccos\big(\ip{a}{b}/(\norm{a}\norm{b} + \varepsilon)\big)$; the code
(lines 58--70) uses the atan2 form above, with the comment that it ``avoids arccos which has
infinite gradient at $\pm 1$''. The code wins. The course (IW24) lists three guards; the code
has a fourth, the $+\varepsilon$ inside each length on lines 55--56 (guard 0 above).
\end{sourcenote}

\begin{worked}[label={we:losses:atan2}]{The atan2 form on pixel A, and on a perfect reconstruction}
\textbf{Pixel A.} $N = 0.147641$, dot $= 0.147$.
\begin{align*}
N^2 &= 0.147641^2 = 0.021798\\
\text{dot}^2 &= 0.147^2 = 0.021609\\
N^2 - \text{dot}^2 &= 0.000189\\
\sqrt{0.000189} &= 0.013748\\
\theta &= \operatorname{atan2}(0.013748,\ 0.147) = \arctan(0.013748/0.147)\\
&= \arctan(0.093524) = 0.09325
\end{align*}
The same angle as arccos gave, with a safe derivative. (The $\varepsilon = 10^{-8}$ under the
root changes the sine term only in the seventh decimal place here.)

\textbf{A perfect reconstruction.} With $\xhat = \vect{x} = (0.1, 0.2, 0.3)$, $N = \norm{\vect{x}}^2
= 0.14$ and $N^2 - \text{dot}^2 = 0$, so the sine term is $\sqrt{0 + 10^{-8}} = 10^{-4}$ and
\[
\theta = \operatorname{atan2}(10^{-4},\ 0.14) = \arctan(0.000714) = 0.000714 \text{ rad}.
\]
The loss never quite reaches zero: guard 3 sets a floor. The real \code{SAMLoss} returns
$0.00084$ here \tagC{} (the floor plus float32 rounding) and a finite gradient; for
$(0.2,0.2,0.2)$ and $(0.25,0.5,0.125)$, the two vectors that broke the naive formula in
\cref{we:losses:nan}, it returns finite angles below $0.001$ rad and finite gradients \tagC.
\end{worked}

\paragraph{A side effect: the floor depends on brightness.} The floor is
$\arctan\!\big(10^{-4}/\norm{\vect{x}}^2\big)$. For a real 165-band pixel at reflectance about
0.2 \tagA{} (an assumed typical level), $\norm{\vect{x}}^2 \approx 165 \times 0.2^2 = 6.6$ and the
floor is $10^{-4}/6.6 = 1.5 \times 10^{-5}$ rad, a thousandth of the 0.02--0.05 rad a good model
reaches. A very dark pixel has a higher floor. This is a tiny breach of scale invariance,
irrelevant to training, and the price of the guard.

\subsection{Masking and averaging the angle}

\begin{lstlisting}
spatial_mask = mask[:, 0:1, :, :]                         # (B, 1, H, W)
num_valid    = spatial_mask.sum().clamp(min=1)            # number of graded pixels, at least 1
sam_loss     = (angles * spatial_mask).sum() / num_valid  # ()  mean angle over graded pixels
\end{lstlisting}
So the spectral-angle term is
\begin{equation}\label{eq:losses:samloss}
L_{\SAM} = \frac{1}{\abs{M}}\sum_{p \in M} \theta_p,
\end{equation}
the mean angle over the graded pixels. Unlike L1 it keeps the spatial shape and divides by an
explicit count, which \code{clamp(min=1)} keeps from being zero.

\subsection{Where arccos is fine}

Two other places compute the spectral angle with plain arccos: the inferencer's
\code{compute_anomaly_scores} and \code{_sam()} in
\file{app/utils/anomaly_detection/scoring.py} (lines 118--128), the numpy version the product
calls. Both run at inference, where no gradient is taken, so the infinite derivative never
arises; only the domain clamp \code{np.clip(cos, -1, 1)} is needed. Numerical-stability
requirements differ between the forward pass and the backward pass. The rounding floor near
$u = 1$ still applies to those scores: an angle below about $2.4 \times 10^{-4}$ rad reads as 0
in float32.

\section{The combined loss and the ramp}\label{sec:losses:ramp}

\subsection{The objective}

\begin{definition}[Combined reconstruction loss]\index{combined loss}
\begin{equation}\label{eq:losses:total}
L(t) = L_1 + \lambda(t)\, L_{\SAM},
\end{equation}
with $L_1$ from \cref{eq:losses:l1}, $L_{\SAM}$ from \cref{eq:losses:samloss}, both over the
same loss mask, and $\lambda(t)$ the \emph{SAM weight} at epoch index $t$
(\code{sam_weight} in the code).
\end{definition}

Its gradient is the sum of the two gradients, the second scaled by $\lambda(t)$, because
differentiation is linear: $\partial L/\partial\hat{x} = \partial L_1/\partial\hat{x} +
\lambda(t)\,\partial L_{\SAM}/\partial\hat{x}$.

\subsection{The ramp}

$\lambda$ is not constant. From
\file{app/foundation_models/trainers/hyperspectral_segformer_mae_trainer.py}, lines 197--204:
\begin{lstlisting}
def _get_sam_weight(self, epoch: int) -> float:
    sam_max     = cfg.sam_weight
    ramp_epochs = cfg.sam_ramp_epochs
    if ramp_epochs <= 0:
        return sam_max
    return sam_max * min(1.0, epoch / ramp_epochs)
\end{lstlisting}
\begin{definition}[Linear ramp]\index{ramp}\index{SAM weight}
With final weight $\lambda_{\max}$ (\code{sam_weight}) and ramp length $R$
(\code{sam_ramp_epochs}),
\begin{equation}\label{eq:losses:ramp}
\lambda(t) = \lambda_{\max}\cdot\min\!\left(1, \frac{t}{R}\right), \qquad t = 0, 1, 2, \dots
\end{equation}
and $\lambda(t) = \lambda_{\max}$ for all $t$ if $R \le 0$.
\end{definition}
A straight line from 0 to $\lambda_{\max}$ over $R$ epochs, then flat forever. Three details
from the source matter.
\begin{itemize}
\item $t$ is the zero-based \emph{epoch index}. The trainer sets \code{self._current_epoch =
epoch} at the start of each epoch (line 53) from \code{range(start_epoch, num_epochs)}, and the
log line prints \code{epoch + 1}. So the logged ``Epoch 1'' trains with $\lambda = 0$, and full
strength arrives at logged epoch $R + 1$.
\item $\lambda$ is updated once per epoch, not per step. Within an epoch it is constant, so the
ramp is really a staircase of $R$ steps (\cref{fig:losses:ramp}).
\item The fine-tuning configurations set \code{sam_ramp_epochs} to 0 (FT02c onwards), so they
run at $\lambda_{\max} = 0.5$ from the first epoch: a model that already reconstructs well does
not need the ramp. (FT01, FT02 and FT02b still had $R = 10$.)\newline
\tagM{research/experiments/indradhanu-finetune/FT04-final-joint/config.json}
\end{itemize}

\subsection{Why ramp at all}

At epoch 0 the parameters are random and the reconstructions have random shape, so the angle
is large, about 0.3--0.4 rad \tagM{configs/hyperspectral_segformer_exp_2.json}, and its gradient
points in a direction that has little to do with the more basic job of getting the levels
roughly right. Optimising both from a random start is unstable; one tends to sabotage the other.
The ramp says: learn the easy thing first. Early, $\lambda \approx 0$ and L1 leads; later, the
shape constraint tightens on a model that already produces sensible magnitudes.

The configuration names the failure mode of turning shape on too fast:
\begin{lstlisting}[language={}]
"L1_increases_when_SAM_ramp_starts": "SAM weight too high (try sam_weight=0.5)"
\end{lstlisting}
If the shape penalty is too strong, the model sacrifices brightness accuracy to satisfy it
and L1 goes \emph{up}.

\subsection{The two shipped configurations}

\begin{center}
\begin{tabular}{lll}
\toprule
setting & v0.1.0 & v0.2.0 \\
\midrule
\code{sam_weight} ($\lambda_{\max}$) & 1.0 & 0.5 \\
\code{sam_ramp_epochs} ($R$) & 20 & 10 \\
\code{mask_ratio} & 0.50 & 0.65 \\
\code{compressed_channels} ($D$) & 24 & 32 \\
\code{drop_rate} & 0.3 & 0.4 \\
\code{learning_rate} & $10^{-4}$ & $10^{-3}$ \\
\code{num_epochs} & 500 & 200 \\
\bottomrule
\end{tabular}
\end{center}
Values from \file{configs/hyperspectral_segformer_exp_1.json} and
\file{configs/hyperspectral_segformer_exp_2.json} \tagM{configs/hyperspectral_segformer_exp_2.json}.
v0.2.0 halved both the final weight and the ramp length: it reaches full strength twice as
fast, but full strength is half as strong.

\begin{worked}[label={we:losses:ramp}]{The v0.2.0 ramp, epoch by epoch}
With $\lambda_{\max} = 0.5$ and $R = 10$, $\lambda(t) = 0.5\min(1, t/10)$.
\begin{center}
\begin{tabular}{rrrrr}
\toprule
epoch index $t$ & logged epoch & $t/10$ & $\min(1, t/10)$ & $\lambda$\\
\midrule
0 & 1 & 0.0 & 0.0 & 0.00\\
1 & 2 & 0.1 & 0.1 & 0.05\\
2 & 3 & 0.2 & 0.2 & 0.10\\
5 & 6 & 0.5 & 0.5 & 0.25\\
9 & 10 & 0.9 & 0.9 & 0.45\\
10 & 11 & 1.0 & 1.0 & 0.50\\
50 & 51 & 5.0 & 1.0 & 0.50\\
199 & 200 & 19.9 & 1.0 & 0.50\\
\bottomrule
\end{tabular}
\end{center}
Row $t = 5$ in full: $5/10 = 0.5$; $\min(1, 0.5) = 0.5$; $0.5 \times 0.5 = 0.25$.
\end{worked}

\begin{figure}[htbp]
\centering
\begin{tikzpicture}
\begin{axis}[width=0.9\linewidth, height=5.2cm, xmin=0, xmax=30, ymin=0, ymax=1.08,
  xlabel={epoch index $t$}, ylabel={SAM weight $\lambda(t)$}, grid=major, grid style={gray!20},
  legend pos=south east, legend style={font=\small}]
\addplot[spec, very thick, const plot] coordinates {(0,0) (1,0.05) (2,0.10) (3,0.15) (4,0.20)
  (5,0.25) (6,0.30) (7,0.35) (8,0.40) (9,0.45) (10,0.5) (30,0.5)};
\addlegendentry{v0.2.0: $\lambda_{\max} = 0.5$, $R = 10$}
\addplot[spat, very thick, dashed, const plot] coordinates {(0,0) (1,0.05) (2,0.10) (3,0.15)
  (4,0.20) (5,0.25) (6,0.30) (7,0.35) (8,0.40) (9,0.45) (10,0.50) (11,0.55) (12,0.60) (13,0.65)
  (14,0.70) (15,0.75) (16,0.80) (17,0.85) (18,0.90) (19,0.95) (20,1.0) (30,1.0)};
\addlegendentry{v0.1.0: $\lambda_{\max} = 1.0$, $R = 20$}
\addplot[loss, thin, domain=0:10] {0.05*x};
\end{axis}
\end{tikzpicture}
\caption{The SAM-weight ramp of both shipped configurations, drawn to scale from
\cref{eq:losses:ramp}. The weight is held for a whole epoch, so each ramp is a staircase; the
thin amber line is the formula $\lambda = 0.05\,t$ that the stairs follow. Both configurations
happen to rise at the same rate, 0.05 per epoch; v0.2.0 stops at half the height. After the
ramp both stay flat for the rest of training (to $t = 199$ or 499).}
\label{fig:losses:ramp}
\end{figure}

\coursefig[0.8\linewidth]{iw-25-ramp}{The course's drawing of the v0.2.0 ramp: the weight on the
angle rises linearly from 0 at epoch index 0 to 0.5 at epoch index 10 and then stays flat, so
L1 leads early training and the angle is phased in.}{fig:losses:courseramp}

\subsection{After the ramp, the angle dominates}

\begin{worked}[label={we:losses:combined}]{The combined loss at the configuration's epoch-100 baseline}
The configuration's baseline at epoch 100 is about $L_1 = 0.010$ and $L_{\SAM} = 0.060$
rad \tagM{configs/hyperspectral_segformer_exp_2.json} (midpoints of its ranges). At $t = 100$,
$\lambda = 0.5$.
\begin{align*}
L &= 0.010 + 0.5 \times 0.060\\
&= 0.010 + 0.030\\
&= 0.040
\end{align*}
Shares: $0.010/0.040 = 25\%$ from L1, $0.030/0.040 = 75\%$ from the angle.
\end{worked}

So after the ramp the shape term carries three times the weight of the brightness term. The
reason is units. L1 is in reflectance and a good value is about 0.01; the angle is in radians
and a good value is about 0.06, six times larger. Multiplying by $\lambda = 0.5$ leaves it three
times larger. $\lambda$ is doing unit conversion as much as weighting: for a hyperspectral model,
spectral shape is meant to be the dominant signal, and it is.

\subsection{Watch the components, not the sum}

During the ramp, $\lambda$ grows, so the combined number can rise while both components fall.
\begin{center}
\begin{tabular}{rrrrl}
\toprule
$t$ & $L_1$ & $L_{\SAM}$ & $\lambda$ & combined\\
\midrule
0 & 0.090 & 0.35 & 0.00 & $0.090 + 0.00 \times 0.35 = 0.090$\\
10 & 0.040 & 0.20 & 0.50 & $0.040 + 0.50 \times 0.20 = 0.040 + 0.100 = 0.140$\\
\bottomrule
\end{tabular}
\end{center}
These are illustrative values inside the configuration's baseline
ranges \tagA. L1 more than halved, the angle nearly halved, and the combined loss rose by
56\% ($0.140/0.090 = 1.56$). The trainer logs the components separately for this reason
(\code{L1: ... | SAM: ... rad (... deg) | lambda: ...}, lines 65--68) and mirrors them to Weights
and Biases as \code{train/l1_loss}, \code{train/sam_loss}, \code{train/sam_loss_deg} and
\code{train/sam_weight}.

\begin{drift}
The configuration's baseline table for v0.1.0's schedule (from
\file{configs/hyperspectral_segformer_exp_2.json}, \code{_loss_reference}; its $\lambda$ column
reflects v0.1.0's 20-epoch ramp to 1.0) \tagM{configs/hyperspectral_segformer_exp_2.json}:

\begin{center}
\begin{tabular}{rlll}
\toprule
epoch & $L_1$ & angle & combined\\
\midrule
0 & 0.08--0.10 & 0.3--0.4 rad (20--25\textdegree) & 0.08--0.10\\
10 & 0.03--0.05 & 0.15--0.25 rad (10--15\textdegree) & 0.10--0.15\\
20 & 0.02--0.03 & 0.08--0.15 rad (5--8\textdegree) & 0.10--0.15\\
100 & 0.005--0.015 & 0.04--0.08 rad (2--5\textdegree) & 0.04--0.08\\
500 & 0.003--0.008 & 0.02--0.05 rad (1--3\textdegree) & 0.02--0.05\\
\bottomrule
\end{tabular}
\end{center}
The combined column rises from epoch 0 to epoch 10 and then falls: the ramp artefact, in the
project's own documented baselines. The same file's red flags: L1 not below 0.05 by epoch 10
(learning rate too low, or data loading); angle stuck above 0.3 after epoch 30 (bottleneck $D$
too narrow, try 32); L1 rising when the ramp starts (\code{sam_weight} too high, try 0.5);
training loss falling but validation not (overfitting). Two of these were hit and fixed between
v0.1 and v0.2: $D$ went 24 to 32 and \code{sam_weight} 1.0 to 0.5.

The v0.2.0 checkpoint chosen for the product had validation loss 0.04349, and the v0.1.0
family's best validation L1 was 0.077 at epoch 92 \tagM{research/model_break_down/05_hyperspectral_segformer_mae.md}.
\end{drift}

\subsection{What \texttt{val\_loss} means, and why it cannot be compared across models}

A checkpoint's \code{val_loss} of 0.04349 is $L_1 + 0.5\,L_{\SAM}$ on the validation patches. It
is not an accuracy, not a percentage and not an L1. A thermal model's \code{val_loss} (Chakshu)
is pure L1 on one band in degrees Celsius: different terms, different units. Comparing the two
numbers is meaningless; models are compared on detection scores (\cref{ch:measures}).

\subsection{Validation grades every pixel exactly once}

Training hides a random 65\% of the valid tokens, so the training loss depends on which tokens
happened to be hidden. Validation removes that luck: \code{compute_validation_loss} (lines
286--342) draws one random half-split of the tokens and runs the model twice, hiding each half
once:
\begin{lstlisting}
rand_mask   = (torch.rand(1, N, device=self.device) > 0.5).float()   # (1, N)
pred_mask_1 = token_validity * (1.0 - rand_mask)    # hide where rand_mask = 0
pred_mask_2 = token_validity * rand_mask            # hide the exact complement
x_hat = x_hat_1 * pass1_pixels + x_hat_2 * pass2_pixels   # each pixel from the pass that hid it
\end{lstlisting}
Every valid pixel is then graded once, from a prediction that did not see it, over the eroded
mask alone (no prediction-mask factor). The number is comparable across epochs, which is what
choosing a checkpoint requires.

\begin{sourcenote}[: the ramped weight in validation]
In \file{app/foundation_models/trainers/hyperspectral_segformer_mae_trainer.py} (line 337), the
validation loss also uses the \emph{ramped} weight $\lambda(t)$. That makes validation numbers
from different epochs of the ramp incomparable for the same reason the training sum is: at $t =
0$ the score is pure L1 and looks best of all. For v0.2.0 the ramp ends at $t = 10$, long
before the chosen checkpoint, so it did not bite. In the SpectralEarth trainer, whose validation
score drives early stopping and \code{best.pt}, it did bite, and that trainer now scores
validation at the full weight from the first epoch
(\file{app/foundation_models/trainers/spectral_stem_mae_trainer.py}, lines 337--341). The story
is told in \cref{ch:pretrain}.
\end{sourcenote}

\begin{keypoint}
A loss is a marking scheme: L1 marks levels, the angle marks shape, and $\lambda(t)$ both
schedules when shape starts to count and converts radians into the same currency as
reflectance.
\end{keypoint}

\begin{confusions}
\confusion{Why average bands first and then pixels? Is averaging an average not wrong?}{It
would be wrong if the groups had different sizes. Every pixel has the same 165 bands, so the
mean of per-pixel means equals the overall mean (\cref{sec:losses:l1}). It is done this way so
that a per-pixel value exists for the mask to select.}
\confusion{Does the loss include pixels the model could see?}{No. The pixel prediction mask
excludes them; grading them would reward copying the input.}
\confusion{Is the spectral angle a distance?}{Not in the usual sense. It is an angle, bounded in
$[0, \pi]$, and two spectra that differ enormously in level can have angle 0
(\cref{we:losses:double}).}
\confusion{The derivative of arccos is infinite at 1. Does the angle's gradient really blow up?}{In
exact arithmetic, no: the gradient has length $1/\norm{\xhat}$ at every nonzero angle
(\cref{sec:losses:samgrad}). The trouble is computing it as a product of an exploding factor
and a vanishing factor in float32, where $u$ near 1 is rounding noise (\cref{we:losses:nan}).}
\confusion{Would normalising each spectrum and then taking L1 give the same thing as the
angle?}{Close, but not identical, and you would lose the brightness information entirely.
Keeping two terms lets you weight them separately.}
\confusion{Why radians and not degrees?}{Radians are what the trigonometric functions produce.
The trainer converts to degrees only when it logs.}
\confusion{Is a lower combined loss always better?}{Only between epochs with the same $\lambda$.
During the ramp the weighting changes, so the combined number is not comparable epoch to
epoch.}
\confusion{Is the L1 loss the same as the L1 anomaly score?}{Same formula, different job. Here it
is averaged into one number to train with; at scoring time it is kept per pixel as a map
(\cref{ch:inference}).}
\end{confusions}

\begin{checkyourself}
\item\label{ex:losses:1} By hand: compute the masked L1 for two graded pixels with two bands
each, $\vect{x}_A = (0.10, 0.20)$, $\xhat_A = (0.13, 0.16)$ and $\vect{x}_B = (0.30, 0.05)$,
$\xhat_B = (0.30, 0.09)$, and give the gradient $\partial L_1/\partial\xhat$ for all four
entries.
\item\label{ex:losses:2} By hand: compute the spectral angle between $\vect{x} = (1, 0)$ and
$\xhat = (1, 1)$ in radians and degrees, then compute $\partial\theta/\partial\xhat$ from
\cref{eq:losses:dtheta}, and check that it is perpendicular to $\xhat$ and has length
$1/\norm{\xhat}$.
\item\label{ex:losses:3} Explain why the derivative of arccos near $u = 1$ matters
\emph{more} as training succeeds, and why, nonetheless, the true gradient of the angle does not
blow up. What does go wrong in float32, and which guard in \code{SAMLoss} prevents it?
\item\label{ex:losses:4} For v0.2.0 ($\lambda_{\max} = 0.5$, $R = 10$): compute $\lambda$ at
epoch index 4, and the combined loss at epoch index 50 if $L_1 = 0.02$ and $L_{\SAM} = 0.12$.
Which term contributes more, and by what factor?
\item\label{ex:losses:5} The five observations $(0.10, 0.11, 0.12, 0.13, 0.90)$ of
\cref{we:losses:median}: compute $S_1(c)$ and $S_2(c)$ at $c = 0.12$ and at $c = 0.272$, and
say which $c$ each loss prefers.
\item\label{ex:losses:6} Why must the token-to-pixel upsampling of the prediction mask be
nearest-neighbour? What exactly goes wrong with bilinear?
\end{checkyourself}

\nextchapter{ch:trainloop}
\par\endgroup
